\section{Theorem Proving}

\subsection{Dependent type theory}
\label{sec:coc}

Examples in this section are based on the excellent introduction to the calculus of inductive constructions by \citet{coc}.

Foundational theorem proving is proving theorems with only the axioms (foundations) of mathematics assumed. Likewise, foundational software verification is
proving theorems (properties) of programs.

Here we focus on the theory of foundational theorem proving and in the next (\ref{sec:meta}) we focus on how these ideas scale to large projects in practice.
Note that we are interested in ``proving theorems'' as software verification is just a particular instance of theorem proving (using program logics, which
is covered in detail in the rest of this report).



We now focus on a specific ``Foundation of Mathematics'' which has proved ideal for theorem proving in practice, as displayed by the most
successful theorem provers - Rocq, Lean, Agda, Idris - all being based on dependent type theory. Dependent type theory has the remarkable property of encoding both
propositions (synonymous to theorems) and proofs on one side and types and programs on the other, as part of the same language. These two sides are in fact formally indistinguishable.
The Curry-Howard isomorphism (or propositions as types method) is an observation that the rules governing logic (introduction elimination rules), are much the same as the inference
rules governing typing judgements. So the two can be used interchangeably. Certain types in our type theory will be interpreted as propositions, and terms of that type will be
proofs of that proposition!

Below is a description of Calculus of Constructions (CoC), a flavour of dependent type theory on which Rocq (CoC and its successors were developed for Rocq) and Lean are based.
Agda is based on Martin Lof Type Theory (MLTT) which differs in a matter of predicativity. We will explain and revisit this briefly at the end.

Pure CoC consists of very few rules and constructs and it is very surprising that it is a valid foundation for all mathematics. We start with a typed lambda calculus in which both the
terms and the types are given by the \emph{same} syntax. This is known as a pure type system (PTS). Instead of function/arrow type we have the dependent product, $\Pi (x : A), B(x)$, where B may be defined using $x$. Additionally,
(as there is no separate syntax for types) all types have \emph{sort}. In other words, every construction in the calculus must have a type, which itself will have a type and so on, infinitely.
$A$ has type $\mathbf{Type_1}$ (or \emph{sort} $\mathbf{Type_1}$, to acknowledge that A is itself a type). There is actually an infinite hierarchy of these types:
$A : \mathbf{Type_1} : \mathbf{Type_2} : \mathbf{Type_3} \ldots$, where the subscripts are indices known as \emph{universes}.

The type of the identity function on A is given by $\Pi (\_ : A), A$. We use the notation $A \to A$ in this case since it is not dependent. For completeness the identity function is typed
$(\lambda x : A \mapsto x) : \Pi (\_ : A), A$
As a first example of a dependent type take the
type of the polymorphic identity function, which takes a type and gives the identity function for that type: $\Pi (A : \mathbf{Type_i}), A \to A$. So $A \to A$ (or $\Pi (\_ : A), A$) is using $A$ in
its definition and is therefore dependent. Again for completeness we type the term which is the polymorphic identity function: 
$\lambda A : \mathbf{Type} \mapsto \lambda x : A \mapsto x : \Pi (A : \mathbf{Type_i}), A \to A$. We will not go into the detail of universe polymorphism which
strictly speaking should be considered since the index $i$ is unquantifed in the above expressions.

We now move to expressing logic in CoC, but before doing so we introduce the sort $\mathbf{Prop} : \mathbf{Type_1}$, so that the (infinite) set of sorts is
given by $\mathcal{S} := \{\mathbf{Prop}\} \cup \bigcup_{i\in\mathbb{N}}\{\mathbf{Type}_i\}$. $\mathbf{Prop}$ is special and is \emph{impredicative} as we will see in \ref{sec:predicativity}.
We map first order logic propositions to types of sort $\mathbf{Prop}$. For example, if $A : \mathbf{Prop}$ and $B : \mathbf{Prop}$, then logical implication $A \Rightarrow B$ corresponds
to the type $A \to B$. This is justified by the meaning of implication being able to prove $B$ given a proof of $A$. Forall quantification (also) corresponds to dependent product types:
given some $T : \mathbf{Type_i}$ and $P : \mathbf{Prop}$, the logical statement $\forall x : T, P(x)$ corresponds to the type $\Pi (x : T), P(x)$. This also makes sense since to prove a
universally quantified statement we need a proof of $P(x)$ for every $x$, which is what a term of type $\Pi (x : T), P(x)$ would give. We will use the notation $\forall x : T, P$ for
$\Pi (x : T), P(x)$, freely, as is standard in Lean.

Dependent type theory as a foundation is an example of \emph{constructive mathematics} because proofs are always built by assembling the required data as seen in
above example. With that intuition in mind, let's show how some other logical connectives are represented:

As seen in this example, the standard recipe in constructive mathematics is to specify a proposition by what is required to prove it. In the next example we
see representation also arises from what the connective enables us to prove. Take $\mathrm{False} := \forall (P : \mathbf{Prop}).\ P : \mathbf{Prop}$.
First notice that if we have a proof of $\mathrm{False}$, that is a term of type $\forall (P : \mathbf{Prop}).\ P$, we would get the proof of any proposition $P$,
out of that, encoding false entails anything. Also, it is not possible to construct a term of type $\forall (P : \mathbf{Prop}).\ P$, ie, there is no proof of $\mathrm{False}$.
A strange thing that one might notice is that we are quantifying over everything in $\mathbf{Prop}$, in order to obtain something of type $\mathbf{Prop}$.

\subsubsection{Predicativity}
\label{sec:predicativity}
\emph{In this section we continue with more examples of logic using dependent type theory since it comes up in a different context in the definition of Iris,
which took a while to figure out.}

Let's now introduce all the typing rules. Remember $\mathcal{S} := \{\mathbf{Prop}\} \cup \bigcup_{i\in\mathbb{N}}\{\mathbf{Type}_i\}$.

\begin{figure}[htbp]
\centering
\small % smaller overall font

\setlength{\arraycolsep}{0.6em} % reduce column spacing

\[
\begin{array}{cc}
\inferrule* [right=\textsc{Type-Prop}]
  {\Gamma \vdash}
  {\Gamma \vdash \mathbf{Prop} : \mathbf{Type}_1}
&
\inferrule* [right=\textsc{Type-Form}]
  {\Gamma \vdash}
  {\Gamma \vdash \mathbf{Type}_i : \mathbf{Type}_{i+1}}
\\[1.2em]

\inferrule* [right=\textsc{Var}]
  {\Gamma \vdash x : A \in \Gamma}
  {\Gamma \vdash x : A}
&
\inferrule* [right=\textsc{Ctx-Ext}]
  {\Gamma \vdash A : s \\ x \notin \Gamma \\ s \in S}
  {\Gamma, x : A \vdash}
\\[1.2em]

\inferrule* [right=\textsc{Fun-Intro}]
  {\Gamma, x : A \vdash t : B}
  {\Gamma \vdash \mathrm{fun}\,x : A \Rightarrow t : \Pi x : A, B}
&
\inferrule* [right=\textsc{Pi-Prop}]
  {\Gamma, x : A \vdash B : \mathbf{Prop}}
  {\Gamma \vdash \Pi x : A, B : \mathbf{Prop}}
\\[1.2em]

\inferrule* [right=\textsc{Pi-Type}]
  {\Gamma \vdash A : \mathbf{Type}_i \\ \Gamma, x : A \vdash B : \mathbf{Type}_j}
  {\Gamma \vdash \Pi x : A, B : \mathbf{Type}_{\max(i,j)}}
&
\inferrule* [right=\textsc{App}]
  {\Gamma \vdash f : \Pi x : A, B \\ \Gamma \vdash a : A}
  {\Gamma \vdash f\,a : B[a/x]}
\\[1.2em]

\multicolumn{2}{c}{
\inferrule* [right=\textsc{Conv}]
  {\Gamma \vdash t : A \\ \Gamma \vdash B : s \\ A \preceq B}
  {\Gamma \vdash t : B}
}
\end{array}
\]

\caption{Inference rules for CoC \citep{coc}}
\end{figure}

We notice that there are two typing rules for the dependent product: \texttt{PI-PROP} for sort $\mathbf{Prop}$ and \texttt{PI-TYPE} for all other sorts. \texttt{PI-TYPE}
is saying that the dependent product falls into a sort which is the max of the sort of its domain and codomain. While there is also a special case rule
\texttt{PI-PROP}, which says that regardless of the domain, if the codomain is in $\mathbf{Prop}$, then the $\Pi$-type stays in $\mathbf{Prop}$.
We say that the sort $\mathbf{Prop}$ is impredicative, because elements in it are defined by quantifying over $\mathbf{Prop}$ itself! Remembering 
$\mathrm{False} := \Pi (P : \mathbf{Prop}).\ P : \mathbf{Prop}$, we see exactly this. And suppose that we instead were to type this with the rule \texttt{PI-TYPE} and
$\mathbf{Prop} \triangleq \mathbf{Type}_0$, then the
domain, $\mathbf{Prop} : \mathbf{Type}_1$ and codomain $P : \mathbf{Prop}$ would give the $\Pi$-type sort $\mathbf{Type}_1$, making $\mathbf{Prop}$ \emph{predicative}. However, we would like $\mathrm{False} : \mathbf{Prop}$,
so the impredicativity of $\mathbf{Prop}$ is desirable. Let's also recall the definition of forall quantification which also depended on predicativity: $\forall x : T, P$,
we can quantify over values in much larger universes or quantify over all $\mathbf{Prop}$s and still end up with a type in $\mathbf{Prop}$. Impredicativity is powerful but fragile.
It can be shown that if $\mathbf{Type}_1$ is also impredicative, then this type theory would become inconsistent \citep{coquand}.

As a final example lets have a look at $A \wedge B \triangleq \forall (P : \mathbf{Prop}).\ (A \to B \to P) \to P : \mathbf{Prop}$. Or making things a little more explicit:
$\wedge := \lambda (A\ B : \mathbf{Prop}) \mapsto \forall (P : \mathbf{Prop}).\ (A \to B \to P) \to P : \mathbf{Prop} \to \mathbf{Prop} \to \mathbf{Prop}$, where we are first using lambda abstractions and then a $\Pi$-type construction.

\subsubsection{Proof Relevance}
Finally, Proof Relevance separates the type theories of Lean and Rocq. Lean is proof irrelevant, meaning terms of types in Prop are definitionally
equal\footnote{Rocq on the other hand is proof relevant by default. However, after noticing that Lean's proof relevance was indispensable for some projects,
they decided to introduce a sort \texttt{SProp}, which is proof irrelevant}. So
for instance if we have a structure bundling data and a proof of a property like so in Lean:

\begin{leancode}
structure MeasurableFun (α β : Type*) [MeasurableSpace α] [MeasurableSpace β] where
  /-- underlying function -/
  toFun : α → β
  meas: Measurable toFun
\end{leancode}

the equivalence of two \texttt{MeasurableFun}s is determined solely by the equality of the \texttt{toFun}, since \texttt{meas} is a (dependent) proof and will be 
definitionally equal if \texttt{toFun} is equal. In practice this is incredibly convenient, and is applied automatically throughout the codebase to simplify proof
objectives. It was always used in an elegant abstraction in establishing that certain constructions form a KRM, which we will detail later.

The choices made in Lean's type theory to bring about proof irrelevance do have some major theoretical drawbacks such as leading to undecidability of definitional
equality \citep[see chapter 3]{mario}, which however don't seem to matter too much in practice.

% \subsubsection{Calculus of Inductive Constructions}
% Mention here that 1) Inductive families... 2) strict positivity


In this section we hope to have gotten across that theorem proving in a dependent type theory is exactly type checking of terms constructed by the user. This is all that is needed
to motivate the next section \label{sec:meta} on how this is partially automated and scaled in Rocq and Lean (going into more Lean specific detail). We also introduced the notions of
predicativity and proof relevance, both of which were encountered over the course of the project. Interestingly they also cleanly separate 3
of the most significant proof assistants.

\subsection{Metaprogramming}
\label{sec:meta}
We now know that constructing a proof involves assembling many subproofs according to the typing rules to obtain a term of a desired Type (actually Prop).
It is however extremely time-consuming and error-prone to go down to this level of detail when constructing a proof. Pen and paper proofs operate at a much
higher level, and the use of metaprogramming in Proof assistants, is motivated by the desire to remove the mental overhead of mechanised proofs, and move closer
to the style of pen and paper proofs.

``Tactics'' are metaprograms meant for constructing proof terms. A Tactic block is started in Lean using $by$ syntax: in our code if we wanted
a term (proof) of type $A : Prop$, we can write $(by \ldots) : A$ (though the type A, would usually be inferred by lean). So how does the $by$ tactic
make it easier to construct the term? It provides a proof context, of premises and goals, which is actually additional state wrapped in a monad called
\texttt{TacticM}. Inside a tactic block a sequence of \texttt{TacticM} monads are composed, each monad constructed using one of Lean's ecosystem of tactics,
incrementally modifying the proof state until ``all goals are solved''. And at the end of a tactic block, using the information acquired in the monadic state
a term of the desired type should be constructed using the information collected in the TacticM monad.

We have been pretty vague about what this ``proof context'' and ``goals'' actually are so far. In lean there are a few more constructors for terms than in the simply
typed lambda calculus, on top of lambda abstraction and application. Terms are called \texttt{Expr} in lean. We introduce the constructor \lean{mvar : MVarID → Expr}.
For building an expression given a metavariable. This is at the heart of how tactics work. Tactics communicate to each other what has been proven and
what is left to prove using metavariables. We will illustrate this through a simple example reproduced from \citep[chapter 4]{meta}. In the proof of the following
theorem we show the evolution of the proof context, line by line.

\begin{leancode}
example {α} (a : α) (f : α → α) (h : ∀ a, f a = a) : f (f a) = a := by
  apply Eq.trans
  · apply h
  · apply h
\end{leancode}

Metavariables are used to represent unknown terms of known types, which may be holes in larger terms or a goal we want to prove (term of a certain Prop that
we are looking to construct). As soon as the tactic block starts with \texttt{by}, a metavariable \texttt{?m1} is created which is unassigned but has type \lean{f (f a) = a}.
Now we look at the state after the first \texttt{apply}. The types of \texttt{Eq.trans} and \texttt{?m1} can be unified, so the tactic succeeds and we have new proof
state as follows:

To a user it would look like this:
\begin{leancode}
...
⊢ f (f a) = ?b

...
⊢ ?b = a

...
⊢ α
\end{leancode}

But it's helpful to see the state of metavariables explicitly:

\begin{leancode}
?m1 := Eq.trans α (f (f a)) ?b a ?m2 ?m3
?m2 : f (f a) = ?b
?m3 : ?b = a
?b 
\end{leancode}

where
 \lean{Eq.trans : ∀ (α : Sort u), ∀ (a b c : α), a = b → b = c → a = c}.

The old goal \texttt{?m1} has been assigned (albeit with holes in the assignment). Two new goals ?m2 and ?m3 were created and an additional mvar ?b,
which is not a Prop was created as well. After the second \texttt{apply}, which acts on the goal ?m2, we assign ?m2 and also assign ?b through the unification process.

\begin{leancode}
?m1 := Eq.trans α (f (f a)) ?b a ?m2 ?m3
?m2 := h (f a)
?m3 : ?b = a
?b  := f a
\end{leancode}
Only ?m3 needs to be assigned which is taken care of by the final \texttt{apply}. At the end of the tactic block, we obtain the term
\lean{Eq.trans α (f (f a)) (f a) a (h (f a)) (h a)} which is typed as \lean{f (f a) = a}, as desired!
